\documentclass[11pt]{ctexart}
\usepackage[a4paper,top=2cm,bottom=2cm,left=2.2cm,right=2.2cm]{geometry}
\usepackage{fancyvrb}
\usepackage{xcolor}
\usepackage{colortbl}
\usepackage{fontspec}
\usepackage[most]{tcolorbox}
\usepackage{booktabs}
\usepackage{enumitem}
\usepackage{titlesec}
\usepackage{hyperref}

\setmonofont{Consolas}
\setCJKmonofont{SimHei}
\setlength{\parindent}{0pt}
\setlist{nosep,leftmargin=1.4em}

% ── Step5 主题色（与 Step1～4 同族、换为深藏青） ──
\definecolor{deep}{RGB}{30,42,88}        % 主色：深藏青
\definecolor{accent}{RGB}{62,84,160}     % 强调：藏青
\definecolor{keyline}{RGB}{90,130,200}   % 知识点：蓝
\definecolor{keybg}{RGB}{238,244,255}
\definecolor{warnline}{RGB}{200,120,60}  % 警告：橙
\definecolor{warnbg}{RGB}{255,243,235}
\definecolor{okline}{RGB}{26,116,88}     % 成果：绿
\definecolor{okbg}{RGB}{240,252,244}
\definecolor{tipline}{RGB}{110,120,180}  % 提示：青灰
\definecolor{tipbg}{RGB}{240,242,250}
\definecolor{codebg}{RGB}{249,249,252}   % 代码底
\definecolor{rowtint}{RGB}{238,240,248}  % 表头浅藏青

\titleformat{\section}{\large\bfseries\color{deep}}{\thesection}{0.6em}{}
\titleformat{\subsection}{\normalsize\bfseries\color{accent}}{\thesubsection}{0.5em}{}
\titleformat{\part}[display]{\centering\Huge\bfseries\color{deep}}{}{0pt}{\Huge}
\titlespacing{\section}{0pt}{14pt}{6pt}
\titlespacing{\subsection}{0pt}{10pt}{4pt}

\newtcolorbox{keybox}[1]{enhanced, colback=keybg, colframe=keyline,
  boxrule=0.6pt, arc=1.5mm, left=3mm, right=3mm, top=2mm, bottom=2mm,
  title={\bfseries #1}, fonttitle=\small\bfseries, coltitle=white,
  colbacktitle=keyline}
\newtcolorbox{warnbox}[1]{enhanced, colback=warnbg, colframe=warnline,
  boxrule=0.6pt, arc=1.5mm, left=3mm, right=3mm, top=2mm, bottom=2mm,
  title={\bfseries #1}, fonttitle=\small\bfseries, coltitle=white,
  colbacktitle=warnline}
\newtcolorbox{okbox}[1]{enhanced, colback=okbg, colframe=okline,
  boxrule=0.6pt, arc=1.5mm, left=3mm, right=3mm, top=2mm, bottom=2mm,
  title={\bfseries #1}, fonttitle=\small\bfseries, coltitle=white,
  colbacktitle=okline}
\newtcolorbox{tipbox}[1]{enhanced, colback=tipbg, colframe=tipline,
  boxrule=0.5pt, arc=1.5mm, left=3mm, right=3mm, top=1.5mm, bottom=1.5mm,
  title={\bfseries #1}, fonttitle=\small\bfseries, coltitle=white,
  colbacktitle=tipline}

\newcommand{\code}[1]{\texttt{#1}}
\DefineVerbatimEnvironment{Verilog}{Verbatim}{frame=single,fontsize=\small}

\begin{document}

% ==================== 封面 ====================
\begin{titlepage}
\centering
\vspace*{2.2cm}
{\color{deep}\Huge\bfseries 拓展功能集}\\[6pt]
{\color{deep}\Large\bfseries 从零手搓 Step5 教程}\\[14pt]
{\large 2026 FPGA 创新设计赛道（安路）赛题二}\\[6pt]
{\large 视觉方向 · Step5\_expansion}\\[20pt]

\begin{tcolorbox}[enhanced, colback=codebg, colframe=accent, boxrule=0.8pt,
  arc=2mm, width=12.5cm, left=4mm, right=4mm, top=3mm, bottom=3mm]
\centering
{\small 读者：零基础队员（建议先读完 Step1～4 教程）\\
目标：把评分表的进阶/扩展项逐个补齐——多算法、运动检测、画质调节、声光\\
前置要求：完成 Step1～4 的手搓或至少精读\\
预计学习时间：10～14 天（每天 2～3 小时）}
\end{tcolorbox}

\vspace{8mm}
\begin{tcolorbox}[enhanced, colback=tipbg, colframe=tipline, boxrule=0.5pt,
  arc=2mm, width=12.5cm, left=4mm, right=4mm, top=2mm, bottom=2mm]
\centering
{\small\color{deep} \textbf{本教程与你手上的代码是什么关系}\\
Step1～4 打通了"看 - 认 - 动 - 显"；Step5 是\textbf{加分项集中营}——\\
官方边缘检测核接入、帧差运动、亮度对比度、三行 OSD、声光反馈。}
\end{tcolorbox}

\vfill
{\small\color{black!50} fpga-ptz-tracker 教程系列 · 2026-09}
\end{titlepage}

\tableofcontents
\newpage

% ============================================================
\part{从零学习教程}

% ============================================================
\section{你在造什么}

\subsection{一句话}

\textbf{对着评分表编程：基础项做完之后，每一个进阶/扩展条目都是一个}
\textbf{独立小功能，这一步把它们一次做齐。}

\subsection{五件套总览}

\begin{center}
\begin{tabular}{llll}
\toprule
\rowcolor{rowtint}\textbf{功能} & \textbf{对应要求} & \textbf{实现} & \textbf{用户看到} \\
\midrule
多算法切换 & 进阶（1） & \code{edge\_detect\_96} + 帧首 mux & 按键切边缘图 \\
延迟可视化 & 进阶（3） & fid 随包传递 & OSD 的 D 字段 \\
运动检测 & 进阶（4） & \code{motion\_detect} & 橙色包围框 \\
三行 OSD+跑马灯 & 扩展（1） & \code{osd\_track} v2/v3 & 文字层 \\
亮度对比度 & 扩展（4） & \code{video\_adj} & 实时调节 \\
声光反馈 & 演示加分 & \code{led\_buzzer} & 绿灯/蜂鸣 \\
\bottomrule
\end{tabular}
\end{center}

\subsection{做完是什么效果}

按一下按键，画面从"彩色跟踪"切换成"黑白边缘图"（无缝，无撕裂）；
有人在镜头前挥手，目标框外多出一个\textbf{橙色}运动包围框；拧亮度旋钮，
画面明暗实时变化且左上角 B/C 字段跟着变；锁定目标时蜂鸣器短鸣一声、
绿灯常亮。

\begin{okbox}{学完 Step5 你能做到}
\begin{itemize}
  \item 把官方例程核"零适配"接进自己的链路
  \item 实现帧边界无缝的算法切换并解释为什么不能帧中切
  \item 写 Q6 定点的亮度/对比度流水线
  \item 用 1bit/像素的块存储器做帧差运动检测
  \item 说出两处"仿真过了上板翻车"的经典陷阱及防范
\end{itemize}
\end{okbox}

% ============================================================
\section{阶段 0：官方例程当积木（1 天）}

\subsection{edge\_detect\_96 是什么}

官方 Lab2 给了一个 Sobel 3x3 边缘检测核（\code{edge\_detect\_96.v}），
工作在 96bit AXI-Stream 上（一次 4 像素），输出边缘点白色、
非边缘黑色。\textbf{端口和官方视频链路完全兼容}：
tuser/tlast/tready 直通，挂上就能用。

\subsection{接入姿势：并行分支，不是替换}

\begin{Verilog}
96bit 入 ----+---- 直通（显示主线不动）
             +---- detect_top（颜色检测分支, 常开）
             +---- edge_detect_96（边缘分支, 常开）
                        |
             mux <- algo_q（帧首锁存的选择信号）
\end{Verilog}

两个算法分支\textbf{同时运行}，切换只动 mux——不是"停一个启一个"。
这样切换零启动延迟，也避免了"新算法第一帧还没出图"的尴尬。

\begin{keybox}{读官方代码的三个动作}
1）读端口和参数（能配置什么）；2）读时序（延迟几拍、tuser/tlast 怎么
传递）；3）读粒度（它按 4 像素/拍工作，输出 2 个判决再复制成 4 个——
了解这点才能对上拍）。本工程把它当\textbf{黑盒积木}用，零修改。
\end{keybox}

\subsection{练习}

\begin{enumerate}
  \item 画出三分支（直通/检测/边缘）的结构图，标出 mux 位置
  \item 思考：边缘分支和检测分支各占多少逻辑资源？为什么要"常开"而不是按需启动？
\end{enumerate}

% ============================================================
\section{阶段 1：帧边界切换的艺术（2 天）}

\subsection{为什么帧中间切会出事故}

mux 前面是两条\textbf{延迟不同的流水线}：检测链约 9 拍 + 平滑行缓存，
边缘核约 5 拍。如果在帧中间拨动选择信号：

\begin{Verilog}
帧中第 300 拍切换:
  第 0~299 像素: 来自 A 链（当前帧第 0~299 像素）
  第 300~ 像素:  来自 B 链（同一帧, 但 B 链进度不同 -> 时间错位）
  -> 画面左右两半来自"不同时刻", 出现垂直错缝
\end{Verilog}

\subsection{解法：帧首拍锁存}

\begin{Verilog}
if (I_tvalid && I_tuser)      // 只在每帧第一拍采样按键
    algo_q <= I_algo_edge;
// 帧内 mux 恒定 -> 切换永远发生在帧与帧之间
\end{Verilog}

切换延迟最多 1 帧（16.7ms）——人眼无感，画面无撕裂。

\subsection{练习}

\begin{enumerate}
  \item 算一算：两条链延迟差 4 拍时，帧中切换的垂直错缝是几个像素？
  \item 思考：为什么采样条件是 \code{I\_tvalid \&\& I\_tuser} 而不是单独 tuser？
\end{enumerate}

% ============================================================
\section{阶段 2：亮度与对比度（1～2 天）}

\subsection{一条公式}

$$out = clamp\big(((c - 128) \times C \gg 6) + 128 + B\big)$$

三个动作：\textbf{居中}（减 128，让"不增不减"的乘法不动中灰）、
\textbf{缩放}（C 是 Q6 定点，64 表示 1.0 倍）、\textbf{平移}（B 有符号加）。
最后 clamp 到 0..255。

\begin{Verilog}
v = c - 128;                    // 居中
v = (v * C) >>> 6;              // Q6: 64=1.0倍, 192=3.0倍
v = v + 128 + B;                // 回中 + 亮度
clamp 到 [0,255]                // 溢出保护
\end{Verilog}

\subsection{两个设计决定}

\begin{keybox}{决定 1：只放显示路径}
调节挂在 OSD 之前的显示通路上，\textbf{检测分支吃的是原图}——
你把画面调到全黑，跟踪照样工作（检测不受显示参数影响）。
这也是评分表"实时调节\textbf{显示参数}"的本意。
\end{keybox}

\begin{keybox}{决定 2：单通道函数复用}
RGB 三通道完全同构，RTL 写一个 \code{adj\_ch} 函数调用三次——
和 Python 金标准一个函数三处调用的结构完全对应，验证时逐通道等价。
\end{keybox}

\subsection{练习}

\begin{enumerate}
  \item 手算：c=100, C=128（2.0 倍）, B=16 时的输出
  \item 手算：c=128 时任何 C/B 组合下输出是多少？为什么？
  \item 思考：为什么 C 用 Q6（64=1.0）而不是直接存小数？
\end{enumerate}

% ============================================================
\section{阶段 3：帧差运动检测（3 天）}

\subsection{原理：一比特的"不同"}

运动 = 这帧和上帧不一样。对\textbf{二值掩码}（Step2 平滑后的 0/1 图）
做帧差：

\begin{Verilog}
diff = cur_mask ^ prev_mask     // 异或: 不一样的地方为 1
cnt = sum(diff)                 // 变化像素总数
bbox = diff 的包围框            // 变化发生在哪里
valid = (cnt >= MOT_TH)         // 变化够多才算"有运动"
\end{Verilog}

\subsection{代价：一帧的"记忆"}

需要存上一帧完整掩码：1280x720x1bit = 115KB——
块 RAM（ERAM）刚好装下。核心结构是\textbf{同地址读旧写新}：

\begin{Verilog}
addr = row_base + x;            // 当前行列
old  = prev_mem[addr];          // 先读: 上帧的值
prev_mem[addr] <= I_mask;       // 后写: 本帧的值盖进去
\end{Verilog}

和 Step2 平滑的索引行缓存同思想——读写同指针，天然对齐。

\subsection{亚门限也要"报告"}

每帧末锁存 cnt/bbox 并输出 1 拍选通脉冲。\textbf{变化少于门限的帧
照样输出}（valid=0）——下游靠 valid 区分"没运动"和"这帧没算"，
可以放心清掉运动指示，不会出现"橙色框赖着不走"。

\subsection{练习}

\begin{enumerate}
  \item 手算：目标从 (10,10) 挪到 (30,20)，8x8 块，diff 像素数是多少？
    （提示：旧位消失 + 新位出现，两块不重叠）
  \item 思考：为什么对"掩码"做帧差而不是对原始 RGB？（提示：存储量）
  \item 估算：如果存 RGB 全帧要做运动检测，需要多少块 RAM？
\end{enumerate}

% ============================================================
\section{阶段 4：仿真与综合的"不一致"陷阱（2 天）}

这是 Step5 最有价值的部分：\textbf{两个仿真里完全看不出、上板才爆的坑}，
都在运动检测模块里真实抓到过。

\subsection{陷阱 1：越界写，仿真沉默、硬件回绕}

帧尾滚动地址 \code{row\_base += W} 后，下一帧 sof 拍若用\textbf{遗留}
row\_base 算写地址，会越界：

\begin{itemize}
  \item \textbf{仿真}：越界写被仿真器\textbf{沉默忽略}——什么都不发生，
    看不出错误；
  \item \textbf{综合}：地址位宽截断，\textbf{回绕}写进 0 号单元——
    把 (0,0) 的前帧参考改写，行为悄悄不同。
\end{itemize}

\begin{warnbox}{铁律一：帧首拍的地址必须显式覆盖}
sof 拍像素就是 (0,0)，把写地址\textbf{强制为 0}（语义本来就正确），
绝不信赖上一帧滚出来的遗留值。
\end{warnbox}

\subsection{陷阱 2：首帧门控靠 X "碰巧"工作}

首帧还没有"上一帧"可比。初版用 \code{started} 标志（首帧 sof 拍置位）：

\begin{itemize}
  \item \textbf{仿真}：RAM 未初始化读出 X，\code{mask \^{} X = X}，
    判决恒假——\textbf{碰巧}没报运动；
  \item \textbf{综合}：RAM 初值默认 0，首帧 diff = mask\^{}0 = mask，
    \textbf{整帧被误判成运动}，橙色框满天飞。
\end{itemize}

\begin{warnbox}{铁律二：门控绝不能依赖 X 传播}
改用 \code{prev\_ready}（\textbf{首帧末行}才置位），首帧不比对不输出。
X 在仿真里是"薛定谔的守卫"——综合一落地就现原形。
\end{warnbox}

\subsection{练习}

\begin{enumerate}
  \item 解释为什么"仿真过、上板错"的 bug 比双边都错的更危险
  \item 排查你写过的模块：有没有"靠 X 挡"或"靠遗留值"的地方？
\end{enumerate}

% ============================================================
\section{阶段 5：三行 OSD 与声光（1～2 天）}

\subsection{OSD 布局}

\begin{Verbatim}
y=16   TRK0 LOCK X0640Y0360     L1: 算法/组号/锁定态/坐标
y=40   D03 B+16 C08 M1839       L2: 延迟/亮度/对比度/运动计数
y=64   2026 FPGA CONTEST - ANLOGIC ...  L3: 跑马灯(青色)
画面:  绿准星 + 红目标框 + 橙运动框
优先级: L1 > L2 > L3 > 准星 > 目标框 > 运动框 > 底图
\end{Verbatim}

L1 前缀随算法变（TRK/EDG），L2 的 D 字段用"显示帧号减包内采集帧号"
实时算出链路延迟，M 字段直接来自运动检测的 cnt。

\subsection{声光反馈}

\begin{Verilog}
锁定:  绿灯常亮
丢失:  红灯 1Hz 闪
锁定沿: 蜂鸣 200ms (从"丢"变"锁"的那一拍触发)
\end{Verilog}

GPIO 直驱，无任何协议——演示时"听见哔一声"就是最好的锁定提示。

\subsection{练习}

\begin{enumerate}
  \item 画出 L1/L2/L3 的像素区域（x 范围、y 范围）
  \item 思考：为什么"锁定沿"触发蜂鸣而不是"锁定电平"？
\end{enumerate}

% ============================================================
\section{阶段 6：跑通回归（1 天）}

\subsection{16 项一键回归}

\begin{Verbatim}
cd D:\fpga-ptz-tracker\tools
python run_all.py          # 全量 16 个测试
\end{Verbatim}

Step5 新增 \code{tb\_motion\_detect}（第 16 项），场景设计覆盖全部
状态转移：

\begin{center}
\begin{tabular}{lll}
\toprule
\rowcolor{rowtint}\textbf{帧} & \textbf{场景} & \textbf{期望} \\
\midrule
f0 & 静止背景 & 只填缓存, 无输出 \\
f1 & 出现 8x8 块 & cnt=64, bbox(10,10)-(17,17) \\
f2 & 块移动 & cnt=128（旧位消失 + 新位出现） \\
f3 & 块消失 & cnt=64（原位置残影） \\
f4 & 5 个孤立像素 & cnt=5, valid=0（\textbf{亚门限拒报}） \\
\bottomrule
\end{tabular}
\end{center}

\begin{okbox}{Step5 学习清单（全部打勾就过关）}
\begin{itemize}
  \item 能说清"并行分支 + 帧首 mux"的切换机制
  \item 能手算 Q6 亮度对比度的一像素变换
  \item 能解释帧差检测的结构（XOR/累加/门限）与亚门限语义
  \item 能复述两处仿真/综合不一致陷阱及铁律
  \item 跑通 16 项回归，说清 motion 测试 5 帧各验什么
\end{itemize}
\end{okbox}

% ============================================================
% Step5 教程 · Part II：工程实现详解
% 由 step5_tutorial.tex \input 引入（共用同一份导言区与主题色）
% ============================================================

\part{工程实现详解}

\begin{tipbox}{这一部分写什么}
Step5 的实现分两类：4 个新 RTL（边缘核接入、画质调节、运动检测、
跑马灯字库）+ 3 个既有模块改造（detect\_top / osd\_track / vision\_top），
外加 7 轮迭代与两处"仿真过、上板会翻车"的修复。
对应目录：\code{03\_vision/step5\_expansion/}。
\end{tipbox}

% ============================================================
\section{工程总览}

\subsection{文件清单}

\begin{center}
\begin{tabular}{llr l}
\toprule
\rowcolor{rowtint}\textbf{文件} & \textbf{类型} & \textbf{行数} & \textbf{职责} \\
\midrule
\code{edge\_detect\_96.v} & 官方 & 149 & Sobel 边缘核（Lab2，零修改复用） \\
\code{video\_adj.v} & RTL & 47 & 亮度/对比度（1 拍流水） \\
\code{motion\_detect.v} & RTL & 117 & 帧差运动检测 v2（修正版） \\
\code{marquee\_rom\_64.v} & 生成 & 1020 & 跑马灯 63 字符字库（脚本生成） \\
\midrule
\code{detect\_top.v} & 改造 & --- & +fid 计数器 +motion 分支挂载 \\
\code{osd\_track.v} & 改造 & --- & v2 三行文本 + 双框（后升级 v3 双缓冲） \\
\code{vision\_top.v} & 改造 & --- & v2 算法 mux + video\_adj + led\_buzzer \\
\midrule
\code{tb/tb\_motion\_detect.v} & 测试 & 117 & 5 帧场景逐字段比对 \\
\bottomrule
\end{tabular}
\end{center}

\subsection{目标包 v1.0 到 v1.1.0}

包从 68bit 扩到 \textbf{135bit}（IF-VISION v1.1.0），新增字段：
algo(1) + fid(8) + mot\_valid(1) + mot\_cnt(14) + mot bbox(62)。
老字段位置不动、新字段插在高位——osd\_track 解包代码按位段
逐一对齐，保留位 bit0 强制补 0（Step4 事故的沉淀）。

% ============================================================
\section{模块详解：edge\_detect\_96（官方核）}

\subsection{粒度对齐：4 像素进、4 像素出}

官方核按 96bit 流的原生粒度工作：每拍 4 像素，Sobel 在灰度上算出
2 个边缘判决，\textbf{复制}回 4 像素输出。这套粒度与我们整条链路一致，
所以挂上即用：

\begin{Verilog}
edge_detect_96 #(.EDGE_TH(80)) u_edge (...);   // 官方默认阈值
// tuser/tlast/tready 直通, 输出再接一个 ser96to24 降带宽
\end{Verilog}

\begin{keybox}{复用官方代码的三条纪律}
1）\textbf{不改}（改动会破坏与官方链路的兼容性，出问题没参照）；\\
2）\textbf{包一层}（阈值等可调参数从顶层传参，不在官方文件里改默认值）；\\
3）\textbf{单独验证}（官方核在 Lab2 已验证，我们验证的是"接线正确"——
tb\_vision\_top 的整帧比对顺带覆盖了边缘分支的 mux 路径）。
\end{keybox}

% ============================================================
\section{模块详解：video\_adj}

47 行的小模块，三处实现决定值得学：

\subsection{Q6 公式的位精确实现}

\begin{Verilog}
function [7:0] adj_ch(input [7:0] c, input [6:0] bright, input [7:0] contrast);
    reg signed [16:0] v;
    begin
        v = $signed({1'b0, c}) - 17'sd128;         // 居中(有符号!)
        v = (v * $signed({1'b0, contrast})) >>> 6; // Q6 缩放(算术右移)
        v = v + 17'sd128 + $signed({1'b0, bright});// 回中 + 亮度
        if (v < 0)         adj_ch = 8'd0;
        else if (v > 255)  adj_ch = 8'd255;
        else               adj_ch = v[7:0];
    end
endfunction
\end{Verilog}

\code{c - 128} 必须有符号（暗部像素是负偏移），乘法后 \code{>>>6}
与 Python 金标准的 floor 移位逐位一致——Step2 沉淀的
"有符号纪律"第三次落地。

\subsection{挂在哪条路上}

\begin{Verilog}
mux(algo_q) -> video_adj -> osd_track -> 显示
                 ^
      检测分支吃的是 mux 之前的原图, 不受调节影响
\end{Verilog}

注意顺序：调节在 \textbf{OSD 之前}——OSD 文字不受亮度调节影响
（调到全黑时状态行依然清晰可读，演示不会"自残"）。

% ============================================================
\section{模块详解：motion\_detect v2}

\subsection{存储与地址}

\begin{Verilog}
(* ram_style = "block" *) reg prev_mem [0:TOTAL-1];  // 1bit/像素, 115KB
wire [20:0] addr = I_sof ? 21'd0 : (row_base + {10'd0, x});
wire        old  = prev_mem[addr];
wire        diff = I_mask ^ old;
// 同地址读旧写新(与 Step2 平滑行缓存同思想)
\end{Verilog}

\subsection{v2 修正的两处（都有实战故事）}

\begin{warnbox}{修正 1：sof 拍写地址强制 0}
帧尾 \code{row\_base += W} 滚动后，下一帧 sof 拍曾用遗留 row\_base：
\textbf{仿真}越界写沉默忽略（(0,0) 参考停在旧值）；
\textbf{综合}ERAM 地址回绕写 0 号单元。同一个 RTL，两种"物性"。
修复：sof 拍地址强制 0——该拍像素恰是 (0,0)，语义本来就正确。
\end{warnbox}

\begin{warnbox}{修正 2：prev\_ready 门控}
初版 \code{started} 在首帧 sof 拍置位，挡不住首帧后续像素：
仿真靠"RAM 读 X 使 diff 恒假"\textbf{碰巧}不报；综合 RAM 初值 0 会
把首帧误判为全帧运动。修复：改 \code{prev\_ready}（首帧末行置位），
首帧不比对不输出。\textbf{门控靠显式状态位，绝不靠 X。}
\end{warnbox}

\subsection{接口语义（写进 IF-VISION）}

每帧末锁存结果并输出 1 拍的 \code{O\_mot\_valid} 选通脉冲。
\textbf{亚门限帧也更新 cnt/bbox 并出选通}，是否算运动由
\code{valid=(cnt>=MOT\_TH)} 表达——下游 OSD 据此可以\textbf{清掉}
运动指示，不会出现"橙色框赖着不走"。已知约定：(0,0) 像素每帧
只写缓存不比对（sof 拍结构，影响 1/921600，金标准同语义）。

% ============================================================
\section{既有模块的改造}

\subsection{detect\_top：+fid +motion 挂载}

\begin{Verilog}
// 进阶(3): 采集帧号, 每输入帧 sof 自增
reg [7:0] fid;
always @(posedge I_clk or negedge I_rst_n)
    if (!I_rst_n) fid <= 8'd0;
    else if (I_tvalid && I_tuser) fid <= fid + 8'd1;

// 进阶(4): 运动分支, 与 centroid 共享平滑掩码流
motion_detect #(.MOT_TH(64)) u_motion (.*, .I_mask(m_mask), ...);
\end{Verilog}

迭代 1 的教训：motion 块一开始插在 \code{m\_*} wire 声明之前，
隐式声明与显式声明冲突编译错——\textbf{Verilog 隐式 wire 是事故之源}，
新代码统一追加在文件尾部（本项目惯例）。

\subsection{osd\_track v2：三行文本与双框}

L1 加算法前缀（TRK/EDG 由包内 algo 位选择），L2 全新遥测行
（D/B/C/M 四字段），L3 跑马灯（\code{marquee\_rom\_64} 生成字库，
63 字符 x 16 行 case 表）。运动框橙色 \code{FF8000}，
优先级插在目标框之后——它信息量最低，被遮挡时最先让路。

\subsection{vision\_top v2：五合一顶层}

\begin{Verilog}
mux(帧首锁存 algo_q) -> video_adj -> osd_track -> OSD'd vpix 出
     |             ^
det 分支(常开)   edge 分支(常开)
135bit pack = {...}  -> cdc -> osd_track
led_buzzer(pix 域)   -> O_led_g/r, O_buzzer
\end{Verilog}

迭代 5：vision\_top 端口加了 3 个新输入（I\_algo\_edge/I\_bright/
I\_contrast），上层 TB 没传——\textbf{端口变更必须同步所有实例化点}，
run\_all 全量回归的价值就是让这类失配 5 分钟内暴露。

% ============================================================
\section{测试平台设计与回归}

\subsection{tb\_motion\_detect：5 帧状态机全覆盖}

64x48 迷你几何（与 centroid 向量同思路：小图跑全逻辑）。
期望值由金标准函数 \code{frame\_diff\_motion}（位于
\code{golden\_model.py}）逐位生成，打包格式为 15 个十六进制字符
（cnt/valid/bbox 六字段移位拼接）。5 帧场景与期望值见
Part I 阶段 6 的表——注意 f4 的\textbf{亚门限拒报}
（cnt=5，valid=0，但 bbox 照常输出）是单独设计的边界用例。

\subsection{回归现状：18 项全绿（本册时点为 16 项 + 一处待办，该待办已补齐）}

\begin{center}
\begin{tabular}{ll}
\toprule
\rowcolor{rowtint}\textbf{状态} & \textbf{说明} \\
\midrule
18/18 PASS & 含 tb\_motion\_detect（Step5 新增第 16 项；闭环/算法切换补强后共 18 项） \\
tb\_osd\_track & 功能断言 + 整帧逐像素比对全过（v2 渲染器补齐后启用） \\
tb\_vision\_top & 直通/云台断言 + 整帧比对全过（同上） \\
\bottomrule
\end{tabular}
\end{center}

\begin{warnbox}{曾经的诚实待办：金标准渲染器 v2（已闭环）}
RTL 的 OSD 已升级到"三行文本+跑马灯"，但 Python 端
\code{render\_osd()} 还是 v1（单行无跑马灯）——比对自然失败。
工程选择\textbf{先禁用帧比对并在 run\_all.py 里标注 pending}，
而不是让回归长期红着。上板前必须补齐 v2 渲染器并恢复比对；
\textbf{功能断言（文字内容/坐标/云台方向）目前由 TB 内部断言承担}。
这条待办写进了迭代记录与结语，不藏不瞒。
\textbf{2026-10-03 更新：v2 渲染器已补齐，逐像素比对全面恢复，全量 18/18 全绿。}
\end{warnbox}

\subsection{Step5 迭代复盘（7 轮）}

\begin{center}
\begin{tabular}{llp{5.2cm}l}
\toprule
\rowcolor{rowtint}\textbf{轮} & \textbf{现象} & \textbf{根因} & \textbf{类型} \\
\midrule
1 & detect\_top 编译错 & motion 块插在声明前，隐式/显式冲突 & RTL \\
2 & started 双重驱动 & 底部残留占位行 & RTL \\
3 & TB 找不到 dut.txt & 重命名 txt 后层次引用未同步 & TB \\
4 & TB 结构性语法错 & 多轮正则改坏 & TB（备份恢复） \\
5 & 上层 TB 编译失败 & 顶层加端口 TB 未传 & TB \\
6 & STEP5\_RTL 未定义 & 脚本变量顺序 & 脚本 \\
7 & 帧比对失败 & 金标准 v1 未跟上 RTL v2 & 流程（已闭环） \\
\bottomrule
\end{tabular}
\end{center}

\begin{keybox}{迭代 4 的处置值得单独学}
TB 被多轮自动修改改坏后，\textbf{从备份恢复原始版本，只重新应用}
必要的 pack/端口改动——"改坏了就回到最后已知好点，重放最小 diff"，
比在废墟上继续打补丁快一个数量级。备份目录
（\code{fpga-ptz-tracker\_backup\_*}）不是摆设。
\end{keybox}

% ============================================================
\section{上板挂载说明（给上板组）}

\begin{center}
\begin{tabular}{ll}
\toprule
\rowcolor{rowtint}\textbf{按键/输入} & \textbf{功能} \\
\midrule
KEY0..2 / 拨码 & preset\_idx 阈值组（0 红 1 绿 2 蓝...） \\
拨码 & track\_en（0 = 云台回中位） \\
按键 & algo\_edge（1 = 边缘显示，帧边界无缝切换） \\
按键x2 & 亮度 -/+（对比度固定或再加两键） \\
GPIO & O\_led\_g / O\_led\_r / O\_buzzer（声光） \\
\bottomrule
\end{tabular}
\end{center}

上板三步：
\begin{enumerate}
  \item 4 个新 RTL + 3 个改造模块加入 TD 工程（edge\_detect\_96 原样复制，
    \textbf{不许改}）；
  \item vision\_top 替换原有顶层实例化位置，端口对照 IF-VISION v1.1.0；
  \item 先空载跑通显示与 OSD，再接舵机与声光，最后现场标定。
\end{enumerate}

\begin{tipbox}{答辩素材提示}
三行 OSD 是"评分表自证屏"（Step4 讲过）；运动检测演示建议安排
"静止 3 秒 - 挥手 - 静止"的节奏，橙色框出现/消失清晰可辨；
声光反馈让"锁定瞬间"可听——三个感官通道同时确认，演示效果拉满。
\end{tipbox}

% ============================================================
\section{小结：Step5 的可复用产出}

\begin{itemize}
  \item \textbf{4 个新模块 + 3 处改造}，进阶（1）（3）（4）与
    扩展（1）（4）全部落地并有验证；
  \item \textbf{两条上板铁律}：帧首地址显式覆盖、门控不靠 X——
    来源于真实的仿真/综合行为差，写进方法论；
  \item \textbf{16 项回归体系}：新增功能的验证模板（迷你几何 + 位精确
    金标准 + 场景覆盖状态机）；
  \item \textbf{一处诚实标注的待办}：渲染器 v2 补齐后恢复整帧比对，
    上板前完成。
\end{itemize}

\begin{okbox}{五步全部完成}
至此视觉方向软件层收官：看（Step1）- 认（Step2）- 动（Step3）-
显（Step4）- 全（Step5），16 项回归全绿，接口文档 v1.1.0，
五册教程 + 全部迭代记录可交付答辩。
\end{okbox}


% ============================================================
\part*{结语}
\clearpage
\section*{总结与下一步}
\addcontentsline{toc}{section}{总结与下一步}

\subsection*{Step5 学到了什么}
\addcontentsline{toc}{subsection}{Step5 学到了什么}
\begin{itemize}
  \item 官方核的复用姿势（端口兼容、粒度对齐、并行常开）
  \item 帧边界锁存切换：延迟 1 帧换画面无缝
  \item Q6 定点的亮度/对比度流水线（居中-缩放-平移-clamp）
  \item 1bit/像素帧差检测：块存储、同地址读旧写新、亚门限语义
  \item 仿真/综合不一致两陷阱（越界写回绕、X 门控失效）与两条铁律
  \item 三行 OSD 与声光反馈的集成
\end{itemize}

\subsection*{下一步做什么}
\addcontentsline{toc}{subsection}{下一步做什么}
\begin{itemize}
  \item \textbf{上板}：五步模块全部挂进官方工程，TD 综合到 PH1P35，
    板上标定阈值与增益（IF-SERVO-PWM / IF-OSD 的现场流程）
  \item \textbf{遗留项}：金标准渲染器升级到 v2（三行文本+跑马灯），
    恢复 tb\_osd\_track 的整帧逐帧比对（详见 Part II"诚实的待办"）
  \item \textbf{答辩}：五册教程 + 全部迭代记录就是最好的过程性证据
\end{itemize}

\end{document}
